\documentclass[11pt,a4paper,danish]{article}
\usepackage[utf8]{inputenc}
\usepackage[danish]{babel} % danske overskrifter
\usepackage[T1]{fontenc} % fonte (output)
\usepackage{graphicx} % indsættelse af billeder
\usepackage{caption}
\usepackage{subcaption}
\usepackage{wrapfig}
\usepackage{geometry}
\usepackage{titlesec}
\usepackage{fancyhdr}
\usepackage{lastpage}
\usepackage{hyperref}
\usepackage{gensymb}

\usepackage{setspace}
\usepackage{tocloft}
\usepackage{parskip}

\usepackage[nottoc]{tocbibind}

%\usepackage{subfig}
%\usepackage{inconsolata}
%\usepackage{enumitem}

\titleformat{\paragraph}[runin]
{\bfseries\scshape}{\theparagraph}{1em}{}

% BEGIN JAVA CODE SECTION
\usepackage{listings}
\usepackage{color}

\definecolor{dkgreen}{rgb}{0,0.6,0}
\definecolor{gray}{rgb}{0.5,0.5,0.5}
\definecolor{mauve}{rgb}{0.58,0,0.82}

\lstset{frame=none,
  language=Java,
  aboveskip=3mm,
  belowskip=3mm,
  showstringspaces=false,
  columns=flexible,
  basicstyle={\small\ttfamily},
  numbers=none,
  numberstyle=\tiny\color{gray},
  keywordstyle=\color{blue},
  commentstyle=\color{dkgreen},
  stringstyle=\color{mauve},
  breaklines=true,
  breakatwhitespace=true,
  tabsize=2
}
% END JAVA CODE SECTION

% BEGIN C++ code
\definecolor{listinggray}{gray}{0.9}
\definecolor{lbcolor}{rgb}{0.9,0.9,0.9}
\lstset{
backgroundcolor=\color{lbcolor},
    tabsize=4,
%   rulecolor=,
    language=[GNU]C++,
        basicstyle=\scriptsize,
        upquote=true,
        aboveskip={1.5\baselineskip},
        columns=fixed,
        showstringspaces=false,
        extendedchars=false,
        breaklines=true,
        prebreak = \raisebox{0ex}[0ex][0ex]{\ensuremath{\hookleftarrow}},
        frame=single,
        numbers=left,
        showtabs=false,
        showspaces=false,
        showstringspaces=false,
        identifierstyle=\ttfamily,
        keywordstyle=\color[rgb]{0,0,1},
        commentstyle=\color[rgb]{0.026,0.112,0.095},
        stringstyle=\color[rgb]{0.627,0.126,0.941},
        numberstyle=\color[rgb]{0.205, 0.142, 0.73},
%        \lstdefinestyle{C++}{language=C++,style=numbers}’.
}
\lstset{
    backgroundcolor=\color{lbcolor},
    tabsize=4,
  language=C++,
  captionpos=b,
  tabsize=3,
  frame=lines,
  numbers=left,
  numberstyle=\tiny,
  numbersep=5pt,
  breaklines=true,
  showstringspaces=false,
  basicstyle=\footnotesize,
%  identifierstyle=\color{magenta},
  keywordstyle=\color[rgb]{0,0,1},
  commentstyle=\color{Darkgreen},
  stringstyle=\color{red}
  }
%END C++ code

%\setcounter{secnumdepth}{3}

\geometry{
 a4paper,
 total={210mm,297mm},
 left=20mm,
 right=20mm,
 top=40mm,
 bottom=20mm,
 }

\pagestyle{fancy}
\fancyhf{}
\rhead{Delta, \date{\today}}
\lhead{Rudy Alex Kohn}
\chead{Uge 2}
\cfoot{Side \thepage \hspace{1pt} af \pageref{LastPage}}

% TOC som links
\usepackage{hyperref}
\hypersetup{
    colorlinks,
    citecolor=black,
    filecolor=black,
    linkcolor=black,
    urlcolor=blue
}

\title{Stripetykkelsesmåler\\Afsluttende Rapport For Praktikophold}
\author{Rudy Alex Kohn}
\date{8. Juni 2017}

\setlength{\parindent}{0pt}

\begin{document}

\maketitle
\renewcommand{\contentsname}{Indeks}
\renewcommand{\listfigurename}{Figurliste}
\renewcommand{\figurename}{Figur}
\renewcommand\refname{Referencer}
%PWNED! :> nice!

\begin{center}
Rudy Alex Kohn\\
ruak@force.dk\\
s133235@student.dtu.dk
\end{center}

\vspace{25mm}

\begin{center}
\includegraphics[scale=0.3]{Billeder/kunlogo.png}
\includegraphics[scale=0.15]{Billeder/DELTA_1024px.png}
\end{center}

\newpage
\setlength\cftparskip{-2pt}
%\setlength\cftbeforechapskip{0.0pt}\singlespacing
\tableofcontents
%\singlespacing

\listoffigures
\listoftables

%1:  produktspecifikation. forudsætninger
%Performance  tykkelsesmålinh præsisosn +/- 0.1 mm
%Product cist max 25,000,- kr gerne mindre.
%2 teknologier og teknologi valg – begrundelse
%3 forsøgsopstilling – gennemgang.
%4 databehandling  procesdiagram.
%5 resulateter  - hvordan ser det ud.
%6 Videre arbejde- plan for den sidste del af projektet- hvor langt når vi, dokumentation etc.

\section{Introduktion}

Denne rapport afdækker praktikforløbet hos Delta (Force, Hørsholm). Den beskriver de forskellige faser i projektet, samt belyser de problematikker der opstod i forbindelse med at løse opgaven.
Praktikopholdet hos Delta er baseret på et afklaringsprojekt med henblik på at få afdækket om hvorhvidt det er muligt at kunne måle højden af en vejstripe (termoplast). I den forbindelse var der fra starten en del problematik, bl.a. var der på forhånd ikke udlagt hvorledes dette skulle foregå. På den baggrund startede forløbet med at der skulle bruges tre uger på indsamling af information omkring hvilken form for teknologi der skulle benyttes. Valget endte med at stå mellem 3d laserskanning eller computer vision.

Det viste sig dog hurtigt at laserskanning var udelukket grundet de høje omkostninger det ville indebære, og der blev derfor fundet et ældre kamera frem der kunne benyttes til projektet.


\section{Produktets specifikationer}

Krav for at produktet rent faktisk kan blive en realitet.

\begin{itemize}
    \item Processen må samlet set ikke tage særlig lang tid, gerne under 4 sekunder.
    \item Målingens præcision skal være på $\pm 0.1$ mm
    \item Omkostningerne skal holde sig under 25.000,- DKr.
\end{itemize}

\section{Teknologivalg}
Denne sektion beskriver de forskellige delelementer projektet består af. Det er både med hensyn til hardware men hvilke udviklingsværktøj der blev benyttet. For at kunne danne et komplet billede, er der information medtaget omkring nogle af beslutninger og deres bevæggrunde. 

\subsection{Kamera}
Der er valgt at benytte sig af et kamera, der på baggrund af en laser der kaster en linje, kan måle forskel i højden fra hvor laserens linje kastes på en termoplastisk markering i forhold til, hvor den rammer ved siden af markeringen.
Kameraet er monokromatisk for at udnytte lyset bedre og har en opløsning på 5 megapixel.

\subsection{Laser}
Laseren kaster en linje ned på markeringsområdet med en styrke der gør det er muligt at kunne aflæse dens lokalisation. Forskellen på hvor laseren rammer på striben og hvor den ikke gør udgør den højde der er interessant.

\subsection{OpenCV}
Til denne proces bliver Open Computer Vision \cite{OpenCV} benyttet, hvilket er et gratis computer vision system der gør det muligt at lave behandling på dataen for at lokalisere laserlinjen på de billeder kameraet tager.

\subsection{Begrundelse}
Valget af kamera og laser er baseret på prisen, et moderne kamera med en tilstrækkelig kvalitet koster under 6.000,- DKr, og en linjelaser koster mellem 2.000,- DKr. og 3.000,- DKr. En laser scannings enhed koster over 15.000,- Dkr. hvilket vil sætte for meget pres på budgettet.

Open Computer Vision er gratis, og kan benyttes i ubegrænset omfang, også i kommercielle forbindelser.

\subsection{C++}
Udviklingen har indtil videre foregået i C$++$, da dette potentielt åbner op for at benytte sig af softwaren, eller dele af den, på generisk vis på tværs af diverse platforme, da softwaren er skrevet med udgangspunkt i standard template library.
Softwaren kræver en compiler der understøtter C++14, for en komplet liste over understøttede compilere, se \href{http://en.cppreference.com/w/cpp/compiler_support}{her}.

\subsection{Udviklingsværktøj}
Softwaren er udviklet ved brug af Microsoft Visual Studio, Community Edition. Selve koden skulle være 100\% kompatibel med andre udviklingskæder som f.eks. GCC. Det ville dog kræve at der blev sat et projekt op der sørger for de korrekte imports.

\subsection{PvAPI}
Kameraet bliver kontrolleret ved hjælp af PvAPI. PvAPI er ellers blevet erstattet af det nyere API fra Allied Vision, Vimba. Men dette API kunne ikke bruges til projektet da det fejler med at kunne kontrollere elementære ting som f.eks. eksponering osv.
Dog skal Vimba's drivere til kameraet været installeret på systemet.


\section{Forsøgsopstillingen}
Denne sektion omhandler opstillingen der er blevet benyttet til at lave test målinger med.



\section{Billedfelt - Opløsning}

Denne sektion omhandler billedfeltet kontra opløsning, og den problematik der ligger bag.

\subsection{Information}

Som udgangspunkt vil der i den afstand der skal tages et billede fra være kendt. Derfor er det muligt at udregne hvor mange pixels der er pr. mm (eller cm). Denne udregning betyder at hvis alle forvrængninger af billedet vil kunne blive taget med i den endelige udregning for en eventuelt højdeforskel baseret ud fra målingspunkterne.

\subsection{Resultat}

Det er på baggrund af informationerne derfor muligt at kunne lave rimelig korrekte udregninger der gør at hver måling bliver så god som mulig. Effekten af at have et mindre område med flere pixels vil kun være en fordel. Det er derfor en fordel at kunne have et kamera med høj opløsning og kvalitet. Så hvis det er muligt skal komprimeringen helt udelukkes eller være minimal da JPEG artifakter kun vil forringe kvaliteten og derved forværre målings præcisionen.


\section{Processen}

Denne sektion omhandler den process der ligger til grund for de resultater der foreligger.

\subsection{Identificering af striben}
For at kunne måle striben, er det nødvendigt at kunne lokaliserer denne. Søgningen efter striben foregår ved at lysindtaget (eksponeringen) bliver trinvis forøget indtil striben er tilstrækkelig synlig i billedet til at kunne identificere den. Når striben er lokaliseret, bliver dette område brugt som reference til de efterfølgende processer. Striben bliver identificeret ud fra dennes sider, hvor kanterne danner linjer på billedet. For at kunne


\section{Fejlede forsøg}
Undervejs i forløbet er forskellige metoder til stadierne i processen blevet afprøvet, hvoraf alle endte med at blive skrottet. Sektionen beskriver de forskellige procedurer, hvilket problem de skulle løse og hvad ideen bag dem var, og hvorfor de ikke var brugbare i praksis.

\subsection{Diagonal matricer}

\subsubsection{Problemstilling}
Identificering af start og slutpunkter for laseren på striben

\subsubsection{Metode}
Metoden er, at opdele billedet i diagonale matricer for derefter at gennemløbe disse for at identificere dobbeltskæringer af laseren. Selve billedet ville blive opdelt i matricer,  i samme antal, som det dobbelte af billedets højde, på begge leder. Det ville resultere i et net af matricer på tværs i begge retninger, hvor pixel intensitetsforskelle for hver i forhold til dens naboer ville medføre et slags toppunkt. I tilfælde af at to var tilstede og de matchede med naboerne, kunne denne information benyttes til at søge efter intensitets variationer, der i sidste ende ville medføre at en potentiel lokation, der kunne bruges, var fundet.

\subsubsection{Resultat}
Metoden viste sig at være utroligt krævende, både med hensyn til design og regnekraft samt noget utilregnelig. Under perfekte situationer virkede det nogenlunde, og hvis der var nogen diskrepans i laserens position grundet støj eller andre elementer ville den give resultater der ikke kunne bruges. Den blev derfor kasseret i et forholdsvis tidligt stadie da det ikke var rentabelt at gå videre i denne retning.

\subsection{Differentiering}

\subsubsection{Problemstilling}
Lokalisering af laser på stribe kontra jord.

\subsubsection{Metode}
Ved at tage alle punkter i X-aksen, og tage forskellen på dem i stigende orden, ville det være muligt at se om et givent koordinat steg eller faldt i højde.
Denne process udføres to gange, for at få absolutte værdier der indikerer udsving i niveau. Ved at finde de største af disse var forhåbningen af det kunne bruges til at finde ud af hvor den stigning laseren var på striben befandt sig.

\subsubsection{Resultat}
I teorien var det en fantastisk ide, men det forblev også kun i teorien. I praksis viste metoden sig at være ubrugelig, fordi afveksling af laserens intensitet på jorden, f.eks. fra fremmedlegemer såsom småsten eller deres refleksioner herfra gjorde det umuligt, at være sikker på om de punkter, der blev lokaliseret rent faktisk var stribens position, eller noget der var en afvigelse fra et helt andet sted.
På trods af det, var det en ganske effektiv og simpel metode til at identificere niveauvariationer.

\subsection{Histogram}

\subsubsection{Problemstilling}
Lokalisering af stribe

\subsubsection{Metode}
Ved at opdele alle billedets pixels i et intensitets histogram kunne de forskellige intensiteter være med til at identificere hvilke der tilhørte striben og derved lokalisere dens position.

\subsubsection{Nødvendigheder}
Algoritmer til at identificere højde- og lavpunkter ud fra data samt en manuel opbygget histogram datastruktur.

\subsubsection{Resultat}
I teorien skulle metoden være mulig, men her opstod et problem der ikke var forudset. Intensiteten af pixels for
\begin{itemize}
	\item Stribe med og uden laser
	\item Jord med og uden laser
\end{itemize}
befandt sig inden for et så lille område at det praktisk var umuligt at foretage udvælgelse af værdierne for videre identificering. Det fremgik også at der, på trods af udelukkelse af for mørke dele, ikke var mulighed for at garantere hvilke intensiteter der tilhørte hvad. Dette forsagede desværre også at det ikke var muligt at søge efter intensitetsgrænser.
Derudover blev situationen heller ikke forbedret af eksponerings problematikker og støv på kamerachippen osv.
Det havde den effekt at det var stort set umuligt at separere de forskellige elementer da overgangene var utroligt ustabile og kunne blive influeret af hvad som helst. På den lyse side blev der dog udviklet to algoritmer, en til at finde højdepunkter og en til at finde lavpunkter, der begge er justerbare med forskellige parametre.
\newpage
\subsection{Nærheds højdemåling}
\subsubsection{Problemstilling}
Undgå indflydelse fra skæv laser ved måling af højden.

\subsubsection{Metode}
Grundet laserens bueform, blev to små bider af laseren udskåret. Hvis disse faldt inde for samme vinkel ville betydningen af laserens bueform have minimal indflydelse. Desværre er denne metode ikke præcis nok da den ikke tager højde for bl.a. udsving i laserens intensitetsvariationer over tid eller elementer der kunne være betydningsfulde for målingen som f.eks. fremmedlegemer.

\subsubsection{Resultat}
Målinger foretaget ved denne metode viser markant andre værdier, og er temmelig svingene i resultater.


\section{Databehandlingen}



\section{Nuværende resultater}

Denne sektion beskriver de resultater, både de endelige samt de metoder og algoritmer brugt og udviklet under forløbet.

\subsection{}



\section{Videre arbejde}


\section{Kommandolinje Interface}
Softwares understøtter styring via et hertil designet kommandolinje interface. Denne sektion beskriver de forskellige muligheder stillet til rådighed ud fra det.

Fungerer ved at starte programmet med de valgte parametrer, som f.eks.\\\\
\texttt{\textbf{ThicknessGauge.exe -d -{}-glob\_name=camera -{}-show\_windows=0}}
\\\\
For en liste over de muligheder der er til rådighed, kan programmet startes med:\\\\
\texttt{\textbf{ThicknessGauge.exe -{}-help}}\\


\subsection{Kommandoerne}

* denoterer ikke komplet understøttet, eller ikke komplet i softwarens nuværende tilstand.

\begin{table}[h]
	\centering
	\label{kommandolinjekontakter}
	\begin{tabular}{lll}
	-d	& -{}-demo & Kører softwaren i dens regulære tilstand \\
	-c	& -{}-calibrate & Kører kamera kalibrering  \\
	-t*	& -{}-test*  &  Kører test tilstand \\
	-g*  & -{}-save\_glob* & Kører en prædefineret opsætning og gemmer en billedserie.
	\end{tabular}
	\caption{Kommandolinje kontakter}
\end{table}

\begin{table}[h]
	\centering
	\label{kommandolinjeindstilling}
	\begin{tabular}{ll}
		-{}-phase\_two\_exposure* & Tvinger programmet til at bruge en eksponering i fase to. \\
-z* & Tvinger aflæsning af område, til nulmålinger. \\
-{}-glob\_name & Brug gemt billedserie, eller "camera" for live feed. \\
-{}-opencv\_threads & Tvinger OpenCV til at bruge N antal tråde. \\
-{}-calibrate\_output & Filnavn til at gemme kamera kalibrering under. \\
-{}-camera\_settings & Kamera kalibrerings fil der skal bruges under processen. \\
-{}-record\_video* & Opretter en video fil af processen. \\
-{}-show\_windows & Optegner processen visuelt. (Giver ikke nøjagtige målinger.) \\
-{}-frames* & Sætter antal frames der skal benyttes under kalkuleringerne. \\
-{}-test\_suite* & Filnavn for at gemme test sæt under. \\
-{}-version & Udskriver programmets version. \\
-{}-help & Viser komplet hjælp af alle kommandolinje parametrer. \\
	\end{tabular}
	\caption{Kommandolinje: indstillinger}
\end{table}


\bibliographystyle{unsrt}
\bibliography{ref}

\end{document}
